\chapter{Конструкторская часть}

\section{Требования к программному обеспечению}
К программе предъявлен ряд требований. Она должна предоставлять пользователю возможность:
\begin{itemize}[label=---]
	\item получить минимальную сумму пути, а также сам путь, используя вышеперечисленные алгоритмы решения задачи коммивояжёра;
	\item провести параметризацию муравьиного алгоритма;
	\item замерить процессорное время работы реализаций алгоритмов решения задачи коммивояжёра.
\end{itemize}


\section{Структуры данных}

Для реализации алгоритмов были использованы следующие структуры данных:
\begin{itemize}[label=---]
	\item матрица~---~массив массивов целого типа;
	\item размерности матрицы~---~значения целого типа.
\end{itemize}

\section{Алгоритмы}

Ниже представлены алгоритмы решения задачи коммивояжёра полным перебором и с помощью муравьиной колонии.

\includeimage
{tspRec}
{f}
{H}
{0.9\textwidth}
{Алгоритм решения задачи коммивояжёра полным перебором}

\includeimage
{ant}
{f}
{H}
{0.85\textwidth}
{Муравьиный алгоритм решения задачи коммивояжёра}
